The p-power assertion in \wikilink{s43--b009}{the source proposition}
provides an elementwise definition valid without a finite-degree hypothesis.

\begin{definition}\label{purely-inseparable-algebraic-extension}
In characteristic $p>0$, an algebraic extension $K/k$ is
\emph{purely inseparable} if for every $\alpha\in K$ there is an
integer $r\ge0$ with $\alpha^{p^r}\in k$.
In characteristic zero the term means $K=k$.
\end{definition}
\begin{proposition}\label{unique-conjugate-criterion}
For an element $\alpha$ algebraic over $k$, its minimal polynomial
has only one distinct root in an algebraic closure if and only if
$k(\alpha)/k$ is purely inseparable.
\end{proposition}
\begin{proof}\label{unique-conjugate-criterion-proof}
In characteristic zero an irreducible polynomial is separable, so
having one distinct root means degree one.
In characteristic $p$, an irreducible monic polynomial $f$ with
zero derivative has the form $g(X^p)$. The polynomial $g$ is still
irreducible, since a factorization of it would factor $f$.
Repeat while the derivative vanishes. Degrees decrease, so eventually
$f(X)=h(X^{p^r})$ with $h$ irreducible and $h'\ne0$.
An irreducible polynomial with nonzero derivative is separable.
In an algebraic closure the map $x\mapsto x^{p^r}$ is bijective,
so $f$ and $h$ have the same number of distinct roots.
If that number is one, $h$ is linear, and $f(X)=X^{p^r}-c$.
Thus $\alpha^{p^r}\in k$.
Conversely, if $\alpha^{p^r}=c\in k$, the minimal polynomial divides
$X^{p^r}-c=(X-\alpha)^{p^r}$ and has only one distinct root.
\end{proof}
\begin{corollary}\label{purely-inseparable-embeddings}
An algebraic extension is purely inseparable if and only if it has
only one embedding over $k$ into a fixed algebraic closure.
\end{corollary}
\begin{proof}\label{purely-inseparable-embeddings-proof}
In a purely inseparable extension, every element has just one possible
image by the proposition, so embeddings agree. An embedding exists
by \wikilink{s37--b012}{embedding extension}.
Conversely, any other root of an element's minimal polynomial gives
an embedding of its simple field. The same extension argument extends
it to the whole algebraic extension. Uniqueness therefore forces
every element to have only one conjugate, and the proposition applies.
\end{proof}
\begin{remark}\label{inseparable-degree-scope}
For finite extensions this agrees with separable degree one and
\wikilink{s43--b008}{the finite inseparable-degree definition}.
For an infinite extension, use the elementwise or embedding condition,
not a quotient of infinite cardinal degrees. In the derivative
argument one stops when the derivative is nonzero, as in the source correction.
\end{remark}
